\documentclass{article}
\usepackage{amsmath, amsthm, amssymb}

\begin{document}

\section*{Theorem and Proof}

\textbf{Definitions.} Let $k, x, x_{eq} \in \mathbb{R}$. We define the force function $F(x)$ and the energy function $E(x)$ as follows:
\begin{align*}
    F(x) &= -k(x - x_{eq}) \\
    E(x) &= \frac{k}{2}(x - x_{eq})^2
\end{align*}

\textbf{Theorem.} The derivative of the energy function with respect to $x$ is equal to the negative of the force function. That is,
\[
\frac{d}{dx} E(x) = -F(x)
\]

\begin{proof}
    We begin by expanding the definition of the energy function on the left-hand side:
    \[
    \frac{d}{dx} E(x) = \frac{d}{dx} \left[ \frac{k}{2}(x - x_{eq})^2 \right]
    \]
    By the constant multiple rule for differentiation, we can factor out the constant $\frac{k}{2}$:
    \[
    \frac{d}{dx} E(x) = \frac{k}{2} \cdot \frac{d}{dx} \left[ (x - x_{eq})^2 \right]
    \]
    Next, we apply the chain rule and the power rule to differentiate the squared term. The outer function is $u^2$ and the inner function is $u = x - x_{eq}$. 
    \[
    \frac{d}{dx} \left[ (x - x_{eq})^2 \right] = 2(x - x_{eq}) \cdot \frac{d}{dx}(x - x_{eq})
    \]
    Since the derivative of $x$ is $1$ and the derivative of the constant $x_{eq}$ is $0$, the inner derivative is simply $1$. Substituting this back yields:
    \[
    \frac{d}{dx} E(x) = \frac{k}{2} \cdot 2(x - x_{eq}) \cdot 1
    \]
    Simplifying the algebraic expression by canceling the $2$s, we get:
    \[
    \frac{d}{dx} E(x) = k(x - x_{eq})
    \]
    Now, we unfold the definition of the negative force function, $-F(x)$:
    \[
    -F(x) = -\left( -k(x - x_{eq}) \right) = k(x - x_{eq})
    \]
    Since both $\frac{d}{dx} E(x)$ and $-F(x)$ simplify to the exact same algebraic expression, we conclude that:
    \[
    \frac{d}{dx} E(x) = -F(x)
    \]
    This completes the proof.
\end{proof}

\end{document}